% =====================================================================
\section{Related work}
\label{sec:lit}
% =====================================================================

\paragraph{Clustering as a measurement problem.}
\textcite{Harris1991Clustering} established the facts every later study works
against: trades concentrate on round fractions, and that concentration rises with
volatility and with the price level, falls with market capitalisation and trading
frequency, and rises as the tick grid gets finer. The last of these matters
directly here, because the Tokyo Stock Exchange operates two grids at once and
moved a large part of its market onto the finer one in 2023.
\textcite{ChungVanNessVanNess2004} add the correlation with the bid--ask spread
that makes spread a necessary control rather than an optional one.

\paragraph{Who does the clustering.}
The behavioural reading is that round numbers are cognitive shortcuts.
\textcite{KuoLinZhao2015} give it teeth by showing that the investors who use
round-price limit orders most heavily also perform worst, which is what licenses
reading a round-price order as an uninformed one.
\textcite{BhattacharyaHoldenJacobsen2012} take the same observation in a
tradable direction, documenting buy--sell imbalances immediately around round
numbers. \textcite{DavisVanNessVanNess2014} supply the fact that organises
Ohta's whole argument: clustering lives in \emph{large} trades, not small ones.

\paragraph{Clustering in the book, not only in the tape.}
\textcite{AhnCaiCheung2005} look at the limit order book directly and find it
thicker at round prices. That is the closest antecedent to the depth-based
measure introduced in Section~\ref{sec:measures}, and the difference is one of
data: a study that sees only executed trades observes round-price orders after
they have been taken, whereas a study that sees the standing book observes them
while they are still exposed. \textcite{BoxGriffith2016} is the closest design
sibling to the present paper, pairing buy/sell asymmetry in clustering with the
price impact of round-price trades.

\paragraph{Noise traders and liquidity.}
Whether noise trading helps or harms liquidity is genuinely unsettled, and this
study cannot settle it. In \textcite{Kyle1985} noise traders are the cover under
which informed traders operate, so more of them means a deeper market.
\textcite{PeressSchmidt2020} find the empirical counterpart: when noise traders
are distracted, liquidity deteriorates. Against that,
\textcite{Foucault1999} formalises pick-off risk --- the cost of leaving a limit
order unrevised --- and \textcite{AquilinaBudishONeill2022} measure the speed
race that determines who bears it. If clustering marks orders that are about to
be picked off, it may mark a market that is about to become more expensive to
trade, not less. Section~\ref{sec:liquidity} reports which way the 2024 data
falls, and treats either answer as a result.

\paragraph{Reading price impact carefully.}
Ohta's second contribution is a warning about interpretation.
\textcite{GlostenMilgrom1985} and \textcite{Hasbrouck1991} make price impact the
canonical measure of information asymmetry, and
\textcite{HendershottJonesMenkveld2011} give the decomposition of the effective
spread into impact and realized spread that this study uses. But in a limit order
market, impact also rises when stale orders sit in the book waiting to be taken,
which has nothing to do with anyone being informed. The impact premium at round
prices reported in Section~\ref{sec:stylized} is exactly that second channel.

\paragraph{Beyond the best quote.}
\textcite{BiaisHillionSpatt1995} is the canonical anatomy of a pure limit order
market of the kind the TSE runs. \textcite{CaoHanschWang2009} justify looking
past the best quote at all, finding that the levels behind it contribute
materially to price discovery; \textcite{NaesSkjeltorp2006} turn the shape of
those levels into the depth-slope measure used here. The order-flow imbalance
construction is that of \textcite{ContKukanovStoikov2014}.

\paragraph{Inference.}
Every panel regression below carries stock and day fixed effects with standard
errors clustered on both margins, following \textcite{CameronGelbachMiller2011}.
With the same stock appearing on 240 days and the same day covering a thousand
stocks, the alternative is not a slightly optimistic standard error but a badly
wrong one.
